\documentclass[10pt,a4paper]{amsart}
\usepackage{amsmath,amssymb,amsthm,graphicx}
% Retained minimal preamble for the Stacks Project chapter more-algebra.tex (More on Algebra).
% The chapter file pinned-original-source.tex carries only \input{preamble} on its first line, so no
% preamble text lives inside the chapter itself.  The Stacks Project uses one single shared preamble for
% every chapter; the snapshot stages it identically as sources/stacks-chow/preamble.tex and
% sources/stacks-etale-cohomology/preamble.tex.  The block below is copied mechanically, line for line,
% from sources/stacks-chow/preamble.tex lines 14-18 (the Xy-pic setup) and 154-217 (the native theorem
% declarations together with the mathematical macros).  Nothing was added, removed or redefined; the
% wrapper class and the amsmath/amssymb/amsthm/graphicx base packages are supplied by the assembler.
% Chapter: more-algebra.tex -- More on Algebra.
% For commutative diagrams we use Xy-pic
\usepackage[all]{xy}

% We use 2cell for 2-commutative diagrams.
\xyoption{2cell}

% Theorem environments.
%
\theoremstyle{plain}
\newtheorem{theorem}[subsection]{Theorem}
\newtheorem{proposition}[subsection]{Proposition}
\newtheorem{lemma}[subsection]{Lemma}

\theoremstyle{definition}
\newtheorem{definition}[subsection]{Definition}
\newtheorem{example}[subsection]{Example}
\newtheorem{exercise}[subsection]{Exercise}
\newtheorem{situation}[subsection]{Situation}

\theoremstyle{remark}
\newtheorem{remark}[subsection]{Remark}
\newtheorem{remarks}[subsection]{Remarks}

\numberwithin{equation}{subsection}

% Macros
%
\def\lim{\mathop{\mathrm{lim}}\nolimits}
\def\colim{\mathop{\mathrm{colim}}\nolimits}
\def\Spec{\mathop{\mathrm{Spec}}}
\def\Hom{\mathop{\mathrm{Hom}}\nolimits}
\def\Ext{\mathop{\mathrm{Ext}}\nolimits}
\def\SheafHom{\mathop{\mathcal{H}\!\mathit{om}}\nolimits}
\def\SheafExt{\mathop{\mathcal{E}\!\mathit{xt}}\nolimits}
\def\Sch{\mathit{Sch}}
\def\Mor{\mathop{\mathrm{Mor}}\nolimits}
\def\Ob{\mathop{\mathrm{Ob}}\nolimits}
\def\Sh{\mathop{\mathit{Sh}}\nolimits}
\def\NL{\mathop{N\!L}\nolimits}
\def\CH{\mathop{\mathrm{CH}}\nolimits}
\def\proetale{{pro\text{-}\acute{e}tale}}
\def\etale{{\acute{e}tale}}
\def\QCoh{\mathit{QCoh}}
\def\Ker{\mathop{\mathrm{Ker}}}
\def\Im{\mathop{\mathrm{Im}}}
\def\Coker{\mathop{\mathrm{Coker}}}
\def\Coim{\mathop{\mathrm{Coim}}}

% Boxtimes
%
\DeclareMathSymbol{\boxtimes}{\mathbin}{AMSa}{"02}

%
% Macros for moduli stacks/spaces
%
\def\QCohstack{\mathcal{QC}\!\mathit{oh}}
\def\Cohstack{\mathcal{C}\!\mathit{oh}}
\def\Spacesstack{\mathcal{S}\!\mathit{paces}}
\def\Quotfunctor{\mathrm{Quot}}
\def\Hilbfunctor{\mathrm{Hilb}}
\def\Curvesstack{\mathcal{C}\!\mathit{urves}}
\def\Polarizedstack{\mathcal{P}\!\mathit{olarized}}
\def\Complexesstack{\mathcal{C}\!\mathit{omplexes}}
% \Pic is the operator that assigns to X its picard group, usage \Pic(X)
% \Picardstack_{X/B} denotes the Picard stack of X over B
% \Picardfunctor_{X/B} denotes the Picard functor of X over B
\def\Pic{\mathop{\mathrm{Pic}}\nolimits}
\def\Picardstack{\mathcal{P}\!\mathit{ic}}
\def\Picardfunctor{\mathrm{Pic}}
\def\Deformationcategory{\mathcal{D}\!\mathit{ef}}

% --- appended 2026-10-12 (Chief): the project own definitions of the reference/slogan/history
% display environments, copied mechanically from the shared Stacks preamble lines 8-12
% (\usepackage{verbatim} plus the three \newenvironment lines, which the chapter preamble keeps
% verbatim from the project preamble). Nothing added, removed or redefined: these environments
% carry no mathematical content and the project itself typesets them as comments. Without these
% three lines a counted statement whose own text uses them fails with Environment slogan undefined.
\usepackage{verbatim}
\newenvironment{reference}{\comment}{\endcomment}
\newenvironment{slogan}{\comment}{\endcomment}
\newenvironment{history}{\comment}{\endcomment}
\usepackage{cleveref}
\crefname{originallemma}{original lemma}{original lemmas}
\Crefname{originallemma}{Original lemma}{Original lemmas}
\makeatletter
\newlabel{algebra-lemma-characterize-reduced-ring-normal}{{030C}{source}{}{sourcepointer-1}{}}
\newlabel{algebra-lemma-characterize-reduced-ring-normal@cref}{{[originallemma][0][]030C}{[1][0][]source}{}{}{}}
\makeatother
\crefname{originallemma}{original lemma}{original lemmas}
\Crefname{originallemma}{Original lemma}{Original lemmas}
\makeatletter
\newlabel{algebra-lemma-colimit-normal-ring}{{037D}{source}{}{sourcepointer-2}{}}
\newlabel{algebra-lemma-colimit-normal-ring@cref}{{[originallemma][0][]037D}{[1][0][]source}{}{}{}}
\makeatother
\crefname{originallemma}{original lemma}{original lemmas}
\Crefname{originallemma}{Original lemma}{Original lemmas}
\makeatletter
\newlabel{algebra-lemma-flat-going-down}{{00HS}{source}{}{sourcepointer-3}{}}
\newlabel{algebra-lemma-flat-going-down@cref}{{[originallemma][0][]00HS}{[1][0][]source}{}{}{}}
\makeatother
\crefname{originallemma}{original lemma}{original lemmas}
\Crefname{originallemma}{Original lemma}{Original lemmas}
\makeatletter
\newlabel{algebra-lemma-integral-going-up}{{00GU}{source}{}{sourcepointer-4}{}}
\newlabel{algebra-lemma-integral-going-up@cref}{{[originallemma][0][]00GU}{[1][0][]source}{}{}{}}
\makeatother
\crefname{originallemma}{original lemma}{original lemmas}
\Crefname{originallemma}{Original lemma}{Original lemmas}
\makeatletter
\newlabel{algebra-lemma-integral-over-field}{{00GS}{source}{}{sourcepointer-5}{}}
\newlabel{algebra-lemma-integral-over-field@cref}{{[originallemma][0][]00GS}{[1][0][]source}{}{}{}}
\makeatother
\crefname{originallemma}{original lemma}{original lemmas}
\Crefname{originallemma}{Original lemma}{Original lemmas}
\makeatletter
\newlabel{algebra-lemma-normal-goes-up}{{033C}{source}{}{sourcepointer-6}{}}
\newlabel{algebra-lemma-normal-goes-up@cref}{{[originallemma][0][]033C}{[1][0][]source}{}{}{}}
\makeatother
\crefname{originallemma}{original lemma}{original lemmas}
\Crefname{originallemma}{Original lemma}{Original lemmas}
\makeatletter
\newlabel{algebra-lemma-quotient-henselization}{{05WQ}{source}{}{sourcepointer-7}{}}
\newlabel{algebra-lemma-quotient-henselization@cref}{{[originallemma][0][]05WQ}{[1][0][]source}{}{}{}}
\makeatother
\crefname{originallemma}{original lemma}{original lemmas}
\Crefname{originallemma}{Original lemma}{Original lemmas}
\makeatletter
\newlabel{algebra-lemma-ring-with-only-minimal-primes}{{04MG}{source}{}{sourcepointer-8}{}}
\newlabel{algebra-lemma-ring-with-only-minimal-primes@cref}{{[originallemma][0][]04MG}{[1][0][]source}{}{}{}}
\makeatother
\crefname{originallemma}{original lemma}{original lemmas}
\Crefname{originallemma}{Original lemma}{Original lemmas}
\makeatletter
\newlabel{lemma-branches}{{31155}{source}{}{sourcepointer-9}{}}
\newlabel{lemma-branches@cref}{{[originallemma][0][]31155}{[1][0][]source}{}{}{}}
\makeatother
\crefname{originallemma}{original lemma}{original lemmas}
\Crefname{originallemma}{Original lemma}{Original lemmas}
\makeatletter
\newlabel{lemma-integral-over-henselian-pair}{{2753}{source}{}{sourcepointer-10}{}}
\newlabel{lemma-integral-over-henselian-pair@cref}{{[originallemma][0][]2753}{[1][0][]source}{}{}{}}
\makeatother
\crefname{originallemma}{original lemma}{original lemmas}
\Crefname{originallemma}{Original lemma}{Original lemmas}
\makeatletter
\newlabel{lemma-irreducible-henselian-pair-connected}{{2948}{source}{}{sourcepointer-11}{}}
\newlabel{lemma-irreducible-henselian-pair-connected@cref}{{[originallemma][0][]2948}{[1][0][]source}{}{}{}}
\makeatother
\title{Fibres of the strict henselization of the normalization and the minimal primes over the algebraic closure of the residue field}
\author{The Stacks Project Authors}
\date{stacks-project:more-algebra.tex:local-snapshot-20261012; GNU FDL 1.2 or later}
\begin{document}
\maketitle
\setcounter{subsection}{1}
\noindent\textit{Editorial note.} All delivered mathematical text is the author's own, copied from the pinned Stacks Project chapter file more-algebra.tex; no statement or proof step was rewritten, reordered or removed.\par\medskip
\noindent\textit{Editorial note.} The chapter file pinned-original-source.tex carries only \textbackslash{}input\{preamble\} on its first line, so no preamble text lives inside the chapter. The retained file sources/stacks-more-algebra/retained-preamble.tex is a mechanical, line-for-line copy of the project's single shared preamble as staged in the snapshot; nothing was added, removed or redefined and no package was added.\par\medskip
\noindent\textit{Editorial note.} The source URL is the source-level Stacks Project URL, as required when more than one card is produced from one source; the per-claim tag is carried by claim\_key.\par\medskip
\noindent\textit{Editorial note.} The author writes 'small detail omitted' once inside the counted proof span, at the quasi-finiteness step, and states at the start that the proof is 'almost exactly the same as for' the earlier branch lemma before supplying the argument in full; both are reproduced verbatim. The eleven references inside the counted proof span are recorded as non-load-bearing navigation pointers.\par\medskip
\noindent\textit{Allowed standard prerequisites (editorial).}\begin{itemize}
\item Henselization and strict henselization of a local ring, and the residue field of the strict henselization. used at statement lines 31346-31347 and proof lines 31384-31388. Fixes the rings whose spectra are compared and supplies the identification of the residue field with kappa\textasciicircum{}\{sep\}.
\item The integral closure of a reduced local ring in its total ring of fractions as a finite product of normal domains. used at proof lines 31400-31402. Is what makes every local ring of (A')\textasciicircum{}\{sh\} a normal domain and hence gives uniqueness of the minimal prime inside each maximal ideal.
\item Henselian pairs, connectedness of V(q) with V of the maximal ideal, and profiniteness of the fibre over the residue field. used at proof lines 31411-31422. Produces the unique maximal ideal containing a given minimal prime, which is the second half of part (3).
\item Separability of the residue field extension and the bijection between Hom into the algebraic closure and the spectrum of a tensor product of fields. used at proof lines 31391-31396. Is the identification in part (1) of the fibre with Hom\_kappa(kappa', kappa-bar).
\end{itemize}\medskip
\ifdefined\hypertarget\hypertarget{sourcepointer-1}{}\fi
\noindent\textit{Source reference.} Original lemma 030C: focused-seeds/AGP/sources/stacks-snapshot-20261011/tags.txt, source lines 3943--3943. Cross-chapter navigation pointer into the Stacks Project; the standard fact it supplies is applied inside the counted proof span and the argument resting on it stays there.\par\medskip
\ifdefined\hypertarget\hypertarget{sourcepointer-2}{}\fi
\noindent\textit{Source reference.} Original lemma 037D: focused-seeds/AGP/sources/stacks-snapshot-20261011/tags.txt, source lines 4196--4196. Cross-chapter navigation pointer into the Stacks Project; the standard fact it supplies is applied inside the counted proof span and the argument resting on it stays there.\par\medskip
\ifdefined\hypertarget\hypertarget{sourcepointer-3}{}\fi
\noindent\textit{Source reference.} Original lemma 00HS: focused-seeds/AGP/sources/stacks-snapshot-20261011/tags.txt, source lines 654--654. Cross-chapter navigation pointer into the Stacks Project; the standard fact it supplies is applied inside the counted proof span and the argument resting on it stays there.\par\medskip
\ifdefined\hypertarget\hypertarget{sourcepointer-4}{}\fi
\noindent\textit{Source reference.} Original lemma 00GU: focused-seeds/AGP/sources/stacks-snapshot-20261011/tags.txt, source lines 620--620. Cross-chapter navigation pointer into the Stacks Project; the standard fact it supplies is applied inside the counted proof span and the argument resting on it stays there.\par\medskip
\ifdefined\hypertarget\hypertarget{sourcepointer-5}{}\fi
\noindent\textit{Source reference.} Original lemma 00GS: focused-seeds/AGP/sources/stacks-snapshot-20261011/tags.txt, source lines 618--618. Cross-chapter navigation pointer into the Stacks Project; the standard fact it supplies is applied inside the counted proof span and the argument resting on it stays there.\par\medskip
\ifdefined\hypertarget\hypertarget{sourcepointer-6}{}\fi
\noindent\textit{Source reference.} Original lemma 033C: focused-seeds/AGP/sources/stacks-snapshot-20261011/tags.txt, source lines 4051--4051. Cross-chapter navigation pointer into the Stacks Project; the standard fact it supplies is applied inside the counted proof span and the argument resting on it stays there.\par\medskip
\ifdefined\hypertarget\hypertarget{sourcepointer-7}{}\fi
\noindent\textit{Source reference.} Original lemma 05WQ: focused-seeds/AGP/sources/stacks-snapshot-20261011/tags.txt, source lines 7594--7594. Cross-chapter navigation pointer into the Stacks Project; the standard fact it supplies is applied inside the counted proof span and the argument resting on it stays there.\par\medskip
\ifdefined\hypertarget\hypertarget{sourcepointer-8}{}\fi
\noindent\textit{Source reference.} Original lemma 04MG: focused-seeds/AGP/sources/stacks-snapshot-20261011/tags.txt, source lines 6042--6042. Cross-chapter navigation pointer into the Stacks Project; the standard fact it supplies is applied inside the counted proof span and the argument resting on it stays there.\par\medskip
\ifdefined\hypertarget\hypertarget{sourcepointer-9}{}\fi
\noindent\textit{Source reference.} Original lemma 31155: sources/stacks-more-algebra/pinned-original-source.tex, source lines 31155--31155. Same-chapter navigation pointer; the statement it names is standard for the section and the argument that consumes it stays inside the counted proof span.\par\medskip
\ifdefined\hypertarget\hypertarget{sourcepointer-10}{}\fi
\noindent\textit{Source reference.} Original lemma 2753: sources/stacks-more-algebra/pinned-original-source.tex, source lines 2753--2753. Same-chapter navigation pointer; the statement it names is standard for the section and the argument that consumes it stays inside the counted proof span.\par\medskip
\ifdefined\hypertarget\hypertarget{sourcepointer-11}{}\fi
\noindent\textit{Source reference.} Original lemma 2948: sources/stacks-more-algebra/pinned-original-source.tex, source lines 2948--2948. Same-chapter navigation pointer; the statement it names is standard for the section and the argument that consumes it stays inside the counted proof span.\par\medskip
\begin{lemma}
\label{lemma-geometric-branches}
Let $(A, \mathfrak m, \kappa)$ be a local ring.
Assume $A$ has finitely many minimal prime ideals.
Let $A'$ be the integral closure of $A$ in the total ring of fractions
of $A_{red}$. Choose an algebraic closure $\overline{\kappa}$ of
$\kappa$ and denote $\kappa^{sep} \subset \overline{\kappa}$
the separable algebraic closure of $\kappa$.
Let $A^{sh}$ be the strict henselization of $A$
with respect to $\kappa^{sep}$.
Consider the maps
$$
\Spec(A') \xleftarrow{c} \Spec((A')^{sh}) \xrightarrow{e} \Spec(A^{sh})
$$
where $(A')^{sh} = A' \otimes_A A^{sh}$. Then
\begin{enumerate}
\item for $\mathfrak m' \subset A'$ maximal the residue field
$\kappa'$ is algebraic over $\kappa$ and the fibre of $c$
over $\mathfrak m'$ can be canonically identified
with $\Hom_\kappa(\kappa', \overline{\kappa})$,
\item the right arrow is bijective on minimal primes,
\item every minimal prime of $(A')^{sh}$ is contained in a unique
maximal ideal and every maximal ideal contains a unique minimal prime.
\end{enumerate}
\end{lemma}

\begin{proof}
The proof is almost exactly the same as for Lemma \ref{lemma-branches}.
Let $I \subset A$ be the ideal of nilpotents.
We have $(A/I)^{sh} = A^{sh}/IA^{sh}$ by
(Algebra, Lemma \ref{algebra-lemma-quotient-henselization}).
The spectra of $A$, $A^{sh}$, $A'$, and $(A')^h$ are the same
as the spectra of $A/I$, $A^{sh}/IA^{sh}$, $A'$, and
$(A')^{sh} = A' \otimes_{A/I} A^{sh}/IA^{sh}$. Thus
we may replace $A$ by $A_{red} = A/I$ and assume $A$ is reduced.
Then $A \subset A'$ which we will use below without further mention.

\medskip\noindent
Proof of (1). The field extension $\kappa'/\kappa$ is algebraic
because $A'$ is integral over $A$. Since $A'$ is integral over $A$,
we see that $(A')^{sh}$ is integral over $A^{sh}$. By going up
(Algebra, Lemma \ref{algebra-lemma-integral-going-up})
every maximal ideal of $A'$, resp.\ $(A')^{sh}$ lies over
the maximal ideal $\mathfrak m$, resp.\ $\mathfrak m^{sh}$ of
$A$, resp.\ $A^h$. We have
$$
(A')^{sh} \otimes_{A^{sh}} \kappa^{sep} =
A' \otimes_A A^h \otimes_{A^h} \kappa^{sep} =
(A' \otimes_A \kappa) \otimes_{\kappa} \kappa^{sep}
$$
because the residue field of $A^{sh}$ is $\kappa^{sep}$.
Thus the fibre of $c$ over $\mathfrak m'$ is the spectrum
of $\kappa' \otimes_\kappa \kappa^{sep}$.
We conclude (1) is true because there is a bijection
$$
\Hom_\kappa(\kappa', \overline{\kappa}) \to
\Spec(\kappa' \otimes_\kappa \kappa^{sep}),\quad
\sigma \mapsto \Ker(
\sigma \otimes 1 : \kappa' \otimes_\kappa \kappa^{sep} \to \overline{\kappa}
)
$$
We will use below that the displayed ring is integral over a field
hence spectrum of this ring is a profinite space, see
Algebra, Lemmas \ref{algebra-lemma-integral-over-field} and
\ref{algebra-lemma-ring-with-only-minimal-primes}.

\medskip\noindent
Proof of (3). The ring $A'$ is a normal ring and in fact a
finite product of normal domains, see
Algebra, Lemma \ref{algebra-lemma-characterize-reduced-ring-normal}.
Since $A^{sh}$ is a filtered colimit of \'etale $A$-algebras,
$(A')^{sh}$ is filtered colimit of \'etale $A'$-algebras
hence $(A')^{sh}$ is a normal ring by
Algebra, Lemmas \ref{algebra-lemma-normal-goes-up} and
\ref{algebra-lemma-colimit-normal-ring}.
Thus every local ring of $(A')^{sh}$ is a normal domain
and we see that every maximal ideal contains a unique minimal prime.
By Lemma \ref{lemma-integral-over-henselian-pair} applied
to $A^{sh} \to (A')^{sh}$ to see that $((A')^{sh}, \mathfrak m(A')^{sh})$
is a henselian pair. If $\mathfrak q \subset (A')^{sh}$
is a minimal prime (or any prime), then the intersection of
$V(\mathfrak q)$ with $V(\mathfrak m (A')^{sh})$ is connected
by Lemma \ref{lemma-irreducible-henselian-pair-connected}
Since $V(\mathfrak m (A')^{sh}) = \Spec((A')^{sh} \otimes \kappa^{sh})$
is a profinite space by we see there is a unique
maximal ideal containing $\mathfrak q$.

\medskip\noindent
Proof of (2). The minimal primes of $A'$ are exactly the primes
lying over a minimal prime of $A$ (by construction).
Since $A' \to (A')^{sh}$ is flat by going down
(Algebra, Lemma \ref{algebra-lemma-flat-going-down})
every minimal prime of $(A')^{sh}$ lies over a minimal prime
of $A'$. Conversely, any prime of $(A')^{sh}$ lying over a
minimal prime of $A'$ is minimal because $(A')^{sh}$ is a filtered
colimit of \'etale hence quasi-finite algebras over $A'$ (small detail omitted).
We conclude that the minimal primes of $(A')^{sh}$
are exactly the primes which lie over a minimal prime of $A$.
Similarly, the minimal primes of $A^{sh}$ are exactly the
primes lying over minimal primes of $A$.
By construction we have $A' \otimes_A Q(A) = Q(A)$ where $Q(A)$
is the total fraction ring of our reduced local ring $A$.
Of course $Q(A)$ is the finite product of residue fields of
the minimal primes of $A$. It follows that
$$
(A')^{sh} \otimes_A Q(A) =
A^{sh} \otimes_A A' \otimes_A Q(A) = A^{sh} \otimes_A Q(A)
$$
Our discussion above shows the spectrum of the ring
on the left is the set of minimal primes of $(A')^{sh}$
and the spectrum of the ring on the right is the
is the set of minimal primes of $A^{sh}$. This finishes the proof.
\end{proof}

\end{document}
